\documentclass[11pt]{article}
\usepackage[margin=1in]{geometry}
\usepackage{amsmath,amssymb,amsthm,mathtools,mathrsfs}
\usepackage[colorlinks=true]{hyperref}
\newtheorem{theorem}{Theorem}[section]
\newtheorem{proposition}[theorem]{Proposition}
\newtheorem{lemma}[theorem]{Lemma}
\newtheorem{corollary}[theorem]{Corollary}
\theoremstyle{remark}
\newtheorem{remark}[theorem]{Remark}
\title{Swing-by: quantitative descent and a noisy learned-mass endpoint theorem}
\author{Theory continuation for review}
\date{September 17, 2026}
\makeatletter
\@namedef{r@sbnew:global}{{source}{}{}{source-anchor}{}}
\@namedef{r@sbpos:scope}{{source}{}{}{source-anchor}{}}
\@namedef{r@sbsign:scope}{{source}{}{}{source-anchor}{}}
\@namedef{r@sbsign:switching}{{source}{}{}{source-anchor}{}}
\@namedef{r@sbwindow:mass}{{source}{}{}{source-anchor}{}}
\makeatother
\begin{document}
\maketitle
\hypertarget{source-anchor}{Cross-references marked ``source refer to the accompanying manuscript.}
\section{Re-auditing the initialized noisy reversal goal}
\label{sbaudit:scope}
\textbf{Historical audit.} The restricted learned-mass result in Section~\ref{sbsep:scope} below supersedes the missing-later-witness status in this earlier section.

This continuation supersedes the architectural assertion that a general noisy
reversal necessarily requires a second large dynamical theory. The intended
target remains initialized noisy many-neuron specialization followed by
reversal on a nonempty OOD sector. Neither that target nor its impossibility
has been proved. In particular, an unproved sign estimate is not a proof that
its eventual proof must be long. We first remove unnecessary uniformity and
grading requirements, and then prove the decreasing side directly from
isotropic initialization.

\subsection{Exact effective matrices and the missing signed quantity}

Use the projected equations in Section~\ref{sbsign:scope}, and put
$b_i=w_{i,\perp}$. For a fixed actual probe pattern $c_i\in\{0,1\}$,
define
\[
 M_c=\sum_i c_i z_iv_i^\top,\qquad
 N_c=\sum_i c_i b_iv_i^\top.
\]
Every neuron is included. On the corresponding cell, for $\xi\in\mathcal S$,
\begin{equation}\label{sbaudit:cell}
 f_\parallel(\xi)=M_c\xi,\qquad f_\perp(\xi)=N_c\xi.
\end{equation}
On an interval with this pattern fixed, the exact equations give
\begin{align}
 \dot M_c&=\sum_i c_i
 (A_iv_iv_i^\top+z_iz_i^\top A_i-z_iP_i^\top),
 \label{sbaudit:matrix-rate}\\
 \dot N_c&=\sum_i c_i
 \left\{-\mathbb E_n[f_\perp(x)(v_i^\top x)_+]v_i^\top
           +b_iz_i^\top A_i-b_iP_i^\top\right\}.
 \label{sbaudit:ambient-rate}
\end{align}
The training selectors in $A_i$ remain the actual selectors. Even if all
active neurons have one common training selector, the two Gram matrices
$\sum c_i v_iv_i^\top$ and $\sum c_i z_iz_i^\top$ enter the rate; fixing
the probe pattern alone does not supply them from $M_c$.

For a unit probe $\xi=\eta z e_p+sk e_q$, where
$k=\sqrt{1-\eta^2z^2}$ and $s\in\{-1,1\}$, write
\[
 \mathcal A_\xi=\eta^{-1}(f_p(\xi)-\eta z),\qquad
 \mathcal T_\xi=\eta^{-2}f_q(\xi).
\]
These definitions do not assume bounded scaled coefficients. Exactly,
\begin{align}
 E(\xi)&=\frac12+\eta^2\left\{
 \frac12\mathcal A_\xi^2-sk\mathcal T_\xi-\frac12z^2
 \right\}+\frac{\eta^4}{2}\mathcal T_\xi^2
                  +\frac12\|f_\perp(\xi)\|^2,
 \label{sbaudit:exact-error}\\
 \frac{\dot E(\xi)}{\eta^2}
 &=\mathcal A_\xi\dot{\mathcal A}_\xi
    +(\eta^2\mathcal T_\xi-sk)\dot{\mathcal T}_\xi
    +\eta^{-2}\langle f_\perp(\xi),\dot f_\perp(\xi)\rangle.
 \label{sbaudit:exact-rate}
\end{align}
For $M_c=\left(\begin{smallmatrix}a&\eta b\\
\eta c&\eta^2d\end{smallmatrix}\right)$, in the ordered basis $(e_p,e_q)$,
$\mathcal A_\xi=(a-1)z+skb$ and $\mathcal T_\xi=cz+skd$.
There is no derivative remainder in \eqref{sbaudit:exact-rate}.
With moving probe patterns, the same identities hold almost everywhere
with frozen-selector rates, as in Proposition~\ref{sbsign:switching}.

There is also a formulation that does not grade any matrix. Write
$r_x=x-f(x)$ and $r_\xi=\xi-f(\xi)$, including ambient output coordinates,
and let $c_i=\mathbf1_{u_i^\top\xi>0}$. Set
\begin{equation}\label{sbaudit:kernel}
 \mathsf K(\xi,x)=\sum_i\left[
 (u_i^\top\xi)_+(u_i^\top x)_+ I_d
 +c_i\alpha_i(x)(\xi^\top x)w_iw_i^\top\right].
\end{equation}
The exact full gradient-flow equations imply, almost everywhere,
\begin{equation}\label{sbaudit:alignment}
 \dot f(\xi)=\mathbb E_n[\mathsf K(\xi,x)r_x],\qquad
 \dot E(\xi)=-\mathfrak J_\xi,\qquad
 \mathfrak J_\xi:=\mathbb E_n\langle r_\xi,\mathsf K(\xi,x)r_x\rangle.
\end{equation}
Indeed, differentiate each $w_i(u_i^\top\xi)_+$ and substitute the two
parameter equations. The positive semidefinite training Gram kernel
controls training-loss dissipation; it does not give a sign for this
mixed probe--training pairing. This is an exact instantaneous Jacobian
kernel, not a frozen-kernel or lazy-training approximation.

To compare with regulator information, set $p=(x_1)_+$,
$D=\mathbb E_n p^2>0$, and define
\[
 d=\mathbb E_n[pr_x]/D,\quad r_x^\circ=r_x-pd,\quad
 q_\xi(x)=\mathsf K(\xi,x)^\top r_\xi,\quad
 Q_\xi=\mathbb E_n[pq_\xi(x)].
\]
Then the exact orthogonal decomposition is
\begin{equation}\label{sbaudit:regulator-split}
 \mathfrak J_\xi=Q_\xi^\top d+
 \mathbb E_n\left\langle q_\xi-pQ_\xi/D,r_x^\circ\right\rangle.
\end{equation}
Regulator progress controls coordinates of $d$. It does not itself give
a signed bound for $Q_\xi$ or the last, centered pairing. This is a precise
additional quantity, rather than a demand for global grading or small
unassigned mass. Reciprocal holding and the signed complement estimates
can be used to estimate this pairing; their existence does not already
establish its sign.

\subsection{An endpoint alternative to derivative propagation}

\begin{proposition}[Three snapshots give a whole sector with a rebound]
\label{sbaudit:secants}
Let a finite network follow an absolutely continuous parameter trajectory
on $[t_0,t_2]$. Suppose $t_0<t_1<t_2$ and a unit $\xi_0\in\mathcal S$
satisfies
\begin{equation}\label{sbaudit:two-gaps}
 E(\xi_0,t_0)-E(\xi_0,t_1)\ge a>0,\qquad
 E(\xi_0,t_2)-E(\xi_0,t_1)\ge b>0.
\end{equation}
Put $L=\max_{j=0,1,2}\sum_i\|w_i(t_j)\|\|u_i(t_j)\|$.
For every unit $\xi\in\mathcal S$ with
\[
 \|\xi-\xi_0\|<\frac{\min(a,b)}{4(1+L)^2},
\]
the two gaps are at least $a/2$ and $b/2$. Consequently the error has an
interior minimum on $[t_0,t_2]$, and its rebound from that minimum to $t_2$
is at least $b/2$. The derivative is negative on a set of positive measure
in $(t_0,t_1)$ and positive on a set of positive measure in $(t_1,t_2)$.
No activation pattern need persist between snapshots.
\end{proposition}
\begin{proof}
At a snapshot the map $f$ is $L$-Lipschitz and positively homogeneous.
On unit directions its residual norm and residual Lipschitz constant are
both at most $1+L$. Thus
$|E(\xi,t_j)-E(\xi_0,t_j)|\le(1+L)^2\|\xi-\xi_0\|$.
Apply this twice to each gap. Continuity gives an attained minimum,
and both endpoints have strictly larger error than $t_1$. Absolute
continuity and integration of the derivative give the two sign sets.
The minimum need not be unique or strict; no such claim follows from
the endpoint inequalities alone.
\end{proof}

At each snapshot form the actual matrices for $\xi_0$ and put
\[
 H_j=(M_j-I)^\top(M_j-I)+N_j^\top N_j.
\]
The two inequalities to prove in this alternative are exactly
\begin{equation}\label{sbaudit:matrix-secants}
 \tfrac12\xi_0^\top(H_0-H_1)\xi_0\ge a,
 \qquad
 \tfrac12\xi_0^\top(H_2-H_1)\xi_0\ge b.
\end{equation}
Each $M_j,N_j$ includes every neuron. A fixed three-snapshot cell permits
using the same quadratic forms for nearby probes; the Lipschitz argument
above does not even need that restriction.

For the stronger claim of strict monotonicity on two open time intervals,
one may instead prove strict signs of \eqref{sbaudit:alignment} at two
regular times, with negative $\dot E$ first. Here regular means every
training preactivation and the chosen probe's preactivations are nonzero.
The parameter vector field and the probe derivative are then continuous
locally. Strict signs persist on two time intervals and one common open
probe sector. This local argument requires no uniform gate persistence
between those intervals. Regularity at those later witnesses must be
checked; the endpoint proposition alone does not assert it.

\subsection{The decreasing side is available from isotropic initialization}

\begin{theorem}[High-probability initial descent for a fixed OOD probe]
\label{sbaudit:initial-descent}
Fix nonzero empirical training samples in $\mathcal S$, a unit probe
$\xi\in\mathcal S$, and $d\ge2$. Initialize independently
$u_i=w_i=\varepsilon\widehat u_i$ with $\widehat u_i$ uniform on
$\mathbb S^{d-1}$. Put $s_0=h\varepsilon^2$,
$C_x=\mathbb E_n\|x\|^2$, and
\[
 F(c)=\sqrt{1-c^2}+(\pi-\arccos c)c,\qquad
 J_\xi=\frac1\pi\mathbb E_n\left[
 (\xi^\top x)\|x\| F\left(\frac{\xi^\top x}{\|x\|}\right)\right].
\]
For $\ell=\log(1/\delta)$ define
\[
 b_h=C_x\sqrt{\frac{6\ell}{d(d+2)h}}+
                 \frac{4C_x\ell}{3h}.
\]
With probability at least $1-\delta$ over initialization,
\begin{equation}\label{sbaudit:initial-rate}
 \dot E(\xi,0)\le-\frac{s_0}{d}
 \{J_\xi-2db_h-2dC_xs_0(2+s_0)\}.
\end{equation}
In particular, if $J_\xi\ge J_*>0$ and
\begin{equation}\label{sbaudit:initial-budget}
 2db_h+2dC_xs_0(2+s_0)\le J_*/2,
\end{equation}
then $\dot E(\xi,0)\le-s_0J_*/(2d)<0$. Almost surely all the finitely
many initial training and probe margins are nonzero, so the strict
decrease persists on some positive time interval and a nonempty open
sector containing $\xi$.
\end{theorem}
\begin{proof}
At initialization, $\|f_0(x)\|\le s_0\|x\|$ and
$\|\mathsf K_0(\xi,x)\|\le2s_0\|x\|$. Replacing both residuals by
their targets in \eqref{sbaudit:alignment} therefore costs at most
\[
 2C_xs_0^2(2+s_0).
\]
Because $w_i=u_i$, the two kernel contributions to the target pairing
are identical. With
\[
 Y_i=\mathbb E_n[(\xi^\top x)(\widehat u_i^\top\xi)_+
                                      (\widehat u_i^\top x)_+],
\]
the target pairing is $2\varepsilon^2\sum_iY_i$.
Rotational integration gives
\[
 \mathbb E[(\widehat u^\top\xi)_+(\widehat u^\top x)_+]
 =\frac{\|x\|}{2\pi d}F(\xi^\top x/\|x\|),
 \qquad \mathbb EY_i=J_\xi/(2d).
\]
For completeness this integral follows by writing a standard Gaussian
as its independent radius and uniform direction, using
$\mathbb E\|g\|^2=d$, and integrating the product of two positive
linear forms over their angular overlap in the two-dimensional plane.
The uniform angular integral is
$\{\sin\theta+(\pi-\theta)\cos\theta\}/(4\pi)$; the independent
two-dimensional Gaussian squared radius has mean $2$.

Also $|Y_i|\le C_x$. Apply Jensen with weights proportional to
$\|x\|^2$, and use
\[
 \mathbb E[(\widehat u^\top a)^2(\widehat u^\top b)^2]
 =\frac{1+2(a^\top b)^2}{d(d+2)}\le\frac3{d(d+2)}
\]
for unit $a,b$. It follows that
$\mathbb EY_i^2\le3C_x^2/[d(d+2)]$.
The centered summands have absolute bound $2C_x$, so the one-sided
Bernstein inequality gives $h^{-1}\sum_iY_i\ge\mathbb EY_i-b_h$
with the stated probability. Substitution proves
\eqref{sbaudit:initial-rate}. On the probability-one nonzero-margin
event the finite empirical vector field is smooth near initialization;
continuity gives the last assertion.
\end{proof}

\begin{corollary}[An explicit noisy SIM region for the decreasing side]
\label{sbaudit:noisy-descent}
Consider equal cluster weights and write
$x^{(k)}=\mu_ke_k+\zeta^{(k)}$, with $\mu_1>\mu_2>0$.
Fix $0<r\le1/4$ and the OOD probe
$\xi_r=re_1-\sqrt{1-r^2}e_2$. Suppose each noise vector has norm at most
$a$, and set
\[
 \overline C=(\mu_1+a)^2,\qquad
 J_*:=\frac{11\mu_1^2r}{32\pi}-2(2\mu_1+a)a>0.
\]
Sufficient primitive initialization conditions are
\begin{equation}\label{sbaudit:primitive-descent}
 h\ge\max\left\{
 \frac{1536\overline C^2\ell}{J_*^2},
 \frac{64d\overline C\ell}{3J_*}\right\},\qquad
 h\varepsilon^2\le\min\left\{1,\frac{J_*}{24d\overline C}\right\}.
\end{equation}
Under these conditions Theorem~\ref{sbaudit:initial-descent} applies.
For planar isotropic Gaussian cluster noise with standard deviations
at most $\sigma_{\max}$ and $N$ samples in total, one may take
$a=\sigma_{\max}\sqrt{2\log(N/\delta_D)}$ with probability at least
$1-\delta_D$. Independently sampling the initialization gives total
success probability at least $1-\delta_D-\delta$.
This is an explicit nonempty region for initial descent, not yet a common
parameter region for specialization and reversal.
\end{corollary}
\begin{proof}
Write $k=\sqrt{1-r^2}$. The noiseless quantity is
\[
 J^{(0)}_{\xi_r}=\frac1{2\pi}
 \{\mu_1^2r[k+(\pi-\arccos r)r]
       -\mu_2^2k[r-k\arcsin r]\}.
\]
Here $k\ge3/4$ and $r-k\arcsin r\le r(1-k)\le r^3$.
Since $\mu_2<\mu_1$ and $r^2\le1/16$, this gives
$J^{(0)}_{\xi_r}\ge11\mu_1^2r/(32\pi)$.

To bound the noise change without differentiating the angular formula,
use the expectation representation in the preceding proof. Cauchy--Schwarz
and $\mathbb E(\widehat u^\top v)^2=\|v\|^2/d$ give, for any $x,y$,
\[
 2d\left|\mathbb E_{\widehat u}
 \bigl[(\xi^\top x)(\widehat u^\top\xi)_+(\widehat u^\top x)_+
       -(\xi^\top y)(\widehat u^\top\xi)_+(\widehat u^\top y)_+\bigr]
 \right|
 \le2(\|x\|+\|y\|)\|x-y\|.
\]
Average this inequality with $y=\mu_ke_k$ to get $J_{\xi_r}\ge J_*$.
The two width bounds make the two terms of $2db_h$ at most $J_*/8$
each. The initialization cap makes the remaining term in
\eqref{sbaudit:initial-budget} at most $J_*/4$. The Gaussian assertion
is the two-dimensional radial tail followed by a union bound. Finally,
choose sufficiently small positive noise, then finite $h$, then positive
$\varepsilon$ sufficiently small to see nonemptiness for this corollary.
\end{proof}

\begin{corollary}[A directional noise estimate reaches noise-scale probe arcs]
\label{sbaudit:sharp-noise}
Let $0<r_0<r_1\le1/4$, suppose all sample perturbations have norm at most
$a\le\mu_2/4$, and suppose the empirical mean perturbation of cluster 1
has norm at most $b$. Uniformly for $r_0\le r\le r_1$,
\begin{equation}\label{sbaudit:sharp-J}
 J_{\xi_r}\ge J_{\rm sh}:=\frac1{2\pi}
 \left\{\mu_1(\mu_1-a)r_0-\mu_1b-a^2
              -\frac{(\mu_2r_1+a)^3}{\mu_2-a}\right\}.
\end{equation}
Whenever $J_{\rm sh}>0$, it may replace $J_*$ in the initialization
conditions, with $\overline C=(\mu_1+a)^2$.
For $n_1$ independent isotropic Gaussian samples in cluster 1, the mean
event can be imposed at failure probability $\delta_m$ with
\[
 b=\sigma_1\sqrt{2\log(1/\delta_m)/n_1}.
\]
\end{corollary}
\begin{proof}
The function $F$ in Theorem~\ref{sbaudit:initial-descent} is increasing,
with $F(0)=1$, so $cF(c)\ge c$ for every $-1\le c\le1$.
On cluster 1 this gives a lower bound by
$\mathbb E_{1,n}[(\xi_r^\top x)\|x\|]$ for the numerator of its
contribution to $J_{\xi_r}$. Since
$|\|x\|-\mu_1|\le a$ and $|\xi_r^\top x|\le\mu_1r+a$,
this is at least
$\mu_1^2r-\mu_1b-\mu_1ra-a^2$.

On cluster 2, $\xi_r^\top x<0$. Write its angle with $\xi_r$ as
$\pi-\beta$, with $0\le\beta<\pi/2$. Then
\[
 F(-\cos\beta)=\sin\beta-\beta\cos\beta
 \le\sin\beta(1-\cos\beta)\le\sin^3\beta.
\]
Moreover $|\det(\xi_r,x)|\le\mu_2r+a$ and $\|x\|\ge\mu_2-a$.
Its possibly negative numerator is therefore bounded below by
$-(\mu_2r+a)^3/(\mu_2-a)$. Combine the two clusters with weights
$1/2$, and use $r\ge r_0$ in the positive term and $r\le r_1$ in the
negative term. The mean-event assertion follows from the planar
Gaussian radial tail with variance $\sigma_1^2/n_1$.
\end{proof}

Unlike the absolute perturbation bound, \eqref{sbaudit:sharp-J} is useful
when $r_0,r_1$ themselves have the noise scale. For fixed means and fixed
$0<\kappa_0<\kappa_1$, take $r_j=\kappa_j e$,
$\sigma_{\max}=O(e)$, polynomial sample sizes growing to infinity,
and fixed positive failure probabilities. Then
$a=O(e\sqrt{\log(1/e)})$, $b=o(e)$, and
\[
 J_{\rm sh}=\frac{\mu_1^2\kappa_0}{2\pi}e+o(e)>0.
\]
Thus the initial-descent argument need not choose probes far outside the
noise-scale learned angles. This asymptotic compatibility concerns this
module only; it does not prove the remaining specialization inequalities.

\begin{corollary}[Uniform descent permits a later, trajectory-chosen probe]
\label{sbaudit:uniform-descent}
Fix $0<r_0<r_1\le1/4$ and the closed probe arc
$\mathcal I=\{re_1-\sqrt{1-r^2}e_2:r_0\le r\le r_1\}$.
Use a positive uniform lower bound $J_*$ from the absolute-noise corollary
(evaluated at $r_0$) or from \eqref{sbaudit:sharp-J}.
Let
\[
 \rho=\frac{J_*}{32d\overline C},\quad
 G=1+\left\lceil\frac{\arcsin r_1-\arcsin r_0}{\rho}\right\rceil,
 \qquad \ell=\log(G/\delta).
\]
Impose \eqref{sbaudit:primitive-descent} with this $\ell$.
With initialization probability at least $1-\delta$, every probe on
$\mathcal I$ has initial upper right error derivative at most
$-s_0J_*/(4d)$. There exists a common $\tau>0$ on which all these errors
are strictly decreasing. Thus a later probe chosen using the trajectory,
anywhere in this deterministic arc, is covered by the same event.
The positive time $\tau$ is not quantitatively lower bounded here.
\end{corollary}
\begin{proof}
Use $G$ equally spaced angular grid points. Every point on the arc is
within Euclidean distance $\rho$ of a grid point. The Bernstein argument
holds at every grid point simultaneously with probability at least
$1-\delta$. Directly from its definition,
\[
 |Y_i(\xi)-Y_i(\zeta)|\le2C_x\|\xi-\zeta\|
 \quad(\|\xi\|=\|\zeta\|=1).
\]
Consequently the transport cost in $2d h^{-1}\sum_iY_i$ is at most
$4dC_x\rho\le J_*/8$. Grid fluctuations cost at most $J_*/4$ and
the residual correction costs at most $J_*/4$. The remaining margin
is at least $3J_*/8$, stronger than asserted.

At an initial zero probe preactivation the target-pairing formula is
unchanged for every compatible selector: its potentially ambiguous
second kernel term has a factor $w_i^\top\xi=u_i^\top\xi=0$.
The residual correction is uniformly bounded for all those selectors.
Training margins are nonzero almost surely, so the parameter vector
field is continuous for a positive initial interval. The graph of
compatible probe selectors is closed. Compactness of the probe arc
and of the selector cube, together with the strict uniform initial
margin, therefore gives a common positive interval on which all
compatible error slopes remain negative. Integrate the almost-everywhere
slopes to obtain strict decrease. This step provides existence of an
interval, not an explicit initialized learning clock.
\end{proof}

\subsection{A local exit mechanism that avoids a global sign theorem}

\begin{proposition}[One gate exit supplies a later positive witness]
\label{sbaudit:exit-witness}
At a time $t_*$ suppose all training preactivations are nonzero. For a
unit probe $\xi_*$ suppose exactly one probe preactivation vanishes,
say $q_j=u_j^\top\xi_*=0$, all other probe preactivations are nonzero,
and $\dot q_j<0$. Remove neuron $j$ and denote the remaining probe
output and its time derivative at $t_*$ by $F$ and $\dot F$.
If
\begin{equation}\label{sbaudit:exit-charge}
 \underbrace{\langle F-\xi_*,\dot F\rangle}_{\text{signed background rate}}
 +\underbrace{\langle F-\xi_*,w_j\rangle\dot q_j}_{\text{exit contribution}}
 >0,
\end{equation}
then the actual full-network probe error has strictly positive derivative
on an interval immediately before $t_*$, and on a common nonempty open
probe sector after shortening that interval.
If $\xi_*$ lies in the interior of $\mathcal I$ and
Corollary~\ref{sbaudit:uniform-descent} holds, there is an initial decreasing
interval and a later increasing interval on a common open OOD sector.
\end{proposition}
\begin{proof}
Before the transverse exit, $q_j>0$ and
$f(\xi_*)=F(t)+w_j(t)q_j(t)$. The active-side derivative is
\[
 \langle F+w_jq_j-\xi_*,
       \dot F+\dot w_jq_j+w_j\dot q_j\rangle.
\]
Its limit at $t_*$ is exactly \eqref{sbaudit:exit-charge}. The regular
training flow and all other probe contributions are smooth locally.
Continuity gives a positive interval before the exit; choose a compact
subinterval strictly before it and then a small open probe sector.
Intersect with the initial descent sector. The strict positive and
negative intervals cannot overlap, so they have the required ordering.
\end{proof}

In the planar notation $v_j=r_j(\cos\phi_j,\sin\phi_j)$,
$z_j=s_j(\cos\psi_j,\sin\psi_j)$, take
$\xi_*=(\sin\phi_j,-\cos\phi_j)$ and suppose $\dot\phi_j>0$.
Then $\dot q_j=-r_j\dot\phi_j$ and
$z_j^\top\xi_*=s_j\sin(\phi_j-\psi_j)$.
If the background were zero, the exit contribution would therefore be
\begin{equation}\label{sbaudit:angular-exit}
 r_js_j\dot\phi_j\sin(\phi_j-\psi_j).
\end{equation}
In the actual network the background in \eqref{sbaudit:exit-charge}
is retained exactly, including every unassigned neuron. Directional
estimates for that signed expression may be cheaper than bounding the
whole complementary mass. Specialization alone supplies neither an
outward exit with output-angle lag nor this signed background bound;
the proposition is a finite-window route to test, not a reached-exit
claim.

\subsection{What the re-audit does and does not decide}

The initialized decreasing side now has a direct high-probability noisy
many-neuron proof. It does not rely on the two-neuron example. The uniform
version covers a prescribed nonempty arc, with a quantitative derivative
margin, but its common initial time interval is only asserted to be positive.
Its width requirement
includes a $d$-dependent term, and is not a dimension-free width theorem.

The first remaining reversal-specific task is to reach, at a probe
covered by that decreasing-side result, either
\begin{equation}\label{sbaudit:later-witness}
 \mathfrak J_\xi(t_+)<0
\end{equation}
at a regular later time, or the second strict secant inequality in
\eqref{sbaudit:two-gaps}. For example, once $t_1>0$ is in the initial
decreasing interval, $E(\xi,t_2)\ge E(\xi,0)$ is a stronger sufficient
endpoint target, but is not necessary and should not be imposed if the
weaker second secant can be proved. A desired post-specialization
decreasing interval would require its own witness; initial descent is
not automatically descent after specialization.

None of the supplied specialization conclusions explicitly bounds the
signed mixed pairing in \eqref{sbaudit:regulator-split}, supplies this
later witness, or supplies the second secant. Nor do the early feedback
and creation guards automatically extend to a later witness time.
Conversely, no argument here excludes deriving that one-sided quantity
from a short finite-window use of the signed equations.

To distinguish non-implication from impossibility, in the noiseless
planar limit put balanced neurons exactly on the two positive axes.
For arbitrary multiplicities the exact flow gives
$f(x)=m_1(x_1)_+e_1+m_2(x_2)_+e_2$, with
$m_j'=2\lambda_jm_j(1-m_j)$ and $0<m_j(0)<1$.
For every fixed probe its error is nonincreasing. This verifies that
perfect axis concentration, fixed gates, increasing progress and balanced
gradient flow do not, by themselves, force reversal. It is not an
isotropic-initialization counterexample, not a noisy-model counterexample,
and not a counterexample to the full basin-creation hypotheses.

Accordingly, neither O1 nor O2 is closed by this continuation, and O3
has not been justified. The Scope A pivot is not accepted on the basis
of the old grading obstruction. The reduced O2 route remains live:
one probe, one finite window, a later signed witness or second secant,
and the sector extension above. Proving that witness from the same
initialized event and closing specialization's common primitive region
are still required before asserting the original headline theorem.

\section{Quantitative descent and specialization-linked endpoint certificates}
\label{sbend:scope}

The original noisy many-neuron specialization-to-reversal goal remains open.
This section makes the initialized decreasing interval quantitative, isolates
the signed endpoint change needed for a post-specialization rebound, and
audits the angular lag without assuming that it decays monotonically.
All conclusions called endpoint certificates below are conditional on their
displayed reached-state hypotheses. No Scope A pivot is made.

\subsection{A deterministic decreasing interval without a minimum gate margin}

\begin{lemma}[Gate-robust control of the initial target pairing]
\label{sbend:pairing-motion}
Let $s_0=h\varepsilon^2\le1$, $C_x\le\overline C$, and
$\|\mathbb E_nxx^\top\|\le\overline\lambda$. On any interval where
$S(t)\le2s_0$, define
\[
 B(t)=\int_0^t\overline\lambda(1+S(v))\sqrt{S(v)}\,dv.
\]
For a unit planar probe and the exact kernel of
\eqref{sbaudit:kernel}, put
$G_\xi(t)=\mathbb E_n[\xi^\top\mathsf K_t(\xi,x)x]$.
Uniformly in the probe and compatible selectors,
\begin{equation}\label{sbend:pairing-motion-bound}
 |G_\xi(t)-G_\xi(0)|
 \le(36+6\sqrt2)\overline C\,\overline\lambda s_0t
 <45\overline C\,\overline\lambda s_0t.
\end{equation}
Also
\begin{equation}\label{sbend:residual-cost}
 |\dot E(\xi,t)+G_\xi(t)|
 \le2\overline C S(t)^2(2+S(t))
\end{equation}
almost everywhere. These estimates do not require fixed training or probe
gates or a lower bound on their initial margins.
\end{lemma}
\begin{proof}
Use stacked matrices $U,W$ whose columns are the full neuron vectors.
Balance gives $\|U\|_F=\|W\|_F=\sqrt S$. The selected full residual
matrix for each neuron has norm at most
$\overline\lambda(1+S)$: its target part is bounded by
$\overline\lambda$, and Cauchy--Schwarz with
$\|f\|_n\le\sqrt{\overline\lambda}S$ bounds the learned part by
$\overline\lambda S$. Thus both parameter displacements have Frobenius
norm at most $B(t)$, and $\|W-U\|_F\le2B(t)$.

Set
\[
 I_\xi(U)=\mathbb E_n\left[(\xi^\top x)
             \sum_i(u_i^\top\xi)_+(u_i^\top x)_+\right].
\]
ReLU Lipschitzness and Cauchy--Schwarz give
\[
 |I_\xi(U(t))-I_\xi(U(0))|
 \le\overline C(\sqrt{S(t)}+\sqrt{s_0})B(t).
\]
In the second kernel term replace
$(w_i^\top\xi)(w_i^\top x)$ by
$(u_i^\top\xi)(u_i^\top x)$, retaining the actual two selectors.
Their product with the latter expression is exactly the ReLU product,
also at zero preactivations. Therefore
\[
 |G_\xi(t)-2I_\xi(U(t))|
 \le\overline C(\|W\|_F+\|U\|_F)\|W-U\|_F
 \le4\overline C\sqrt{S(t)}B(t).
\]
At initialization the difference is zero. Since
$B(t)\le3\overline\lambda\sqrt{2s_0}t$, combining the two estimates
proves \eqref{sbend:pairing-motion-bound}. Finally
$\|f(x)\|\le S\|x\|$ and
$\|\mathsf K_t(\xi,x)\|\le2S\|x\|$ bound the cost of replacing both
residuals in \eqref{sbaudit:alignment} by their targets, proving
\eqref{sbend:residual-cost}.
\end{proof}

\begin{theorem}[Explicit uniform initial descent time and snapshot drop]
\label{sbend:descent-clock}
On the initialization event proved in
Corollary~\ref{sbaudit:uniform-descent}, with its primitive width and
initialization conditions, define
\begin{equation}\label{sbend:tau-dec}
 \tau_{\rm dec}:=\min\left\{
 \frac{\log(5/4)}{6\overline\lambda},
 \frac{J_*}{720d\overline C\,\overline\lambda}\right\},
 \qquad c_{\rm dec}:=\frac{s_0J_*}{4d}.
\end{equation}
Then every probe on the deterministic arc $\mathcal I$ satisfies
\[
 \dot E(\xi,t)\le-c_{\rm dec}
 \quad\text{for almost every }0\le t\le\tau_{\rm dec}.
\]
In particular, the deterministic time and uniform drop
\begin{equation}\label{sbend:early-drop}
 t_{\rm early}:=\tau_{\rm dec}/2,\qquad
 a_{\rm dec}:=\frac{s_0J_*\tau_{\rm dec}}{8d}>0
\end{equation}
obey $E(\xi,0)-E(\xi,t_{\rm early})\ge a_{\rm dec}$.
The arc's angular width is exactly
\[
 \Delta_{\rm init}=\arcsin r_1-\arcsin r_0.
\]
For $r_j=\kappa_je$, it lies between
$(\kappa_1-\kappa_0)e$ and
$(\kappa_1-\kappa_0)e/\sqrt{1-r_1^2}$.
\end{theorem}
\begin{proof}
Before the residual correction, the uniform-grid event gives
$G_\xi(0)\ge5s_0J_*/(8d)$ for every probe. The grid fluctuation costs
$J_*/4$ and the transport costs $J_*/8$ from a target lower bound $J_*$.
The installed exponential estimate in Lemma~\ref{sbwindow:mass} gives
\[
 S(t)\le\alpha s_0,\qquad \alpha=(5/4)^{1/3}<2
 \quad(t\le\tau_{\rm dec}).
\]
The initial residual cost obeys
$2\overline C s_0^2(2+s_0)\le s_0J_*/(4d)$ by the primitive
initialization cap. Its bound at time $t$ is at most $\alpha^3$ times
this value, hence at most $5s_0J_*/(16d)$.
Lemma~\ref{sbend:pairing-motion} costs at most $s_0J_*/(16d)$ on the
chosen interval. Subtracting these two costs from $5s_0J_*/(8d)$ leaves
$s_0J_*/(4d)$. Integrate the almost-everywhere derivative bound; the
absolute continuity statement includes gate changes. The two arc-width
bounds follow by integrating $(\arcsin r)'=(1-r^2)^{-1/2}$.
\end{proof}

\subsection{The correct time bookkeeping for an O2b endpoint result}

Initial descent supplies a drop at $t_{\rm early}$. It does not supply
$E(0)>E(t_{\rm spec})$. Consequently an initial drop and a distinct
post-specialization rebound naturally use four snapshots, not a tacit
identification of $t_{\rm early}$ with a specialization time.

\begin{proposition}[Four snapshots and an explicit final sector]
\label{sbend:four-snapshots}
Suppose $0<t_{\rm early}<t_{\rm spec}\le t_a<t_b\le H$,
and $\xi_0$ is a probe in the central half of the angular arc
$\mathcal I$. In addition to Theorem~\ref{sbend:descent-clock}, suppose
\[
 E(\xi_0,t_b)-E(\xi_0,t_a)\ge b_*>0.
\]
Put
\begin{align}
 S_H&:=\min\{s_0e^{2\overline\lambda H},
                         s_0+\overline C H/2\},\label{sbend:mass-cap}\\
 \rho_*&:=\min\left\{1,
       \frac{\min(a_{\rm dec},b_*)}{4(1+S_H)^2}\right\},\notag\\
 \Delta_{\rm swing}&:=\min\{\Delta_{\rm init}/2,
                                  4\arcsin(\rho_*/2)\}>0.
 \label{sbend:sector-width}
\end{align}
There is an open sector centered at $\xi_0$ of angular width
$\Delta_{\rm swing}$ on which the initial drop is at least
$a_{\rm dec}/2$ and the post-specialization rebound is at least $b_*/2$.
Each probe error has an interior minimum on $[0,t_b]$ and a rebound from
that minimum of at least $b_*/2$. This is an O2b-endpoint certificate when
the stipulated specialization state is actually certified. It does not
place that minimum after $t_{\rm spec}$.
\end{proposition}
\begin{proof}
The exponential and linear mass bounds, respectively
Lemmas~\ref{sbwindow:mass} and~\ref{sbnew:global}, give $S\le S_H$.
Neuronwise balance implies
$\sum_i\|w_i\|\|u_i\|=S$, so this also bounds the snapshot Lipschitz
constant. Apply the proof of Proposition~\ref{sbaudit:secants} to the
two separate pairs of times. A chord radius $\rho_*$ is an angular
half-width $2\arcsin(\rho_*/2)$. Intersect with the initial arc; the
central-half assumption leaves an angular half-width at least
$\Delta_{\rm init}/4$.
The global minimum is interior because the left endpoint exceeds the
value at $t_{\rm early}$ and the right endpoint exceeds the value at
$t_a$, even though those two intermediate values need not be ordered.
\end{proof}

\subsection{Which transverse output removal can actually cause a rebound?}

For a fixed $\xi=re_1-ke_2$, $k=\sqrt{1-r^2}$, put
\[
 X=f_1(\xi)-r,\qquad \Lambda=-f_2(\xi),\qquad P=f_\perp(\xi).
\]
The exact snapshot identity is
\begin{equation}\label{sbend:leakage-secants}
 E_b-E_a=(\Lambda_a-\Lambda_b)
          \left[k-\frac{\Lambda_a+\Lambda_b}{2}\right]
       +\frac{X_b^2-X_a^2}{2}
       +\frac{\|P_b\|^2-\|P_a\|^2}{2}.
\end{equation}
In active-matrix notation, this scalar is
\[
 \Lambda_j=k(M_j)_{22}-r(M_j)_{21}.
\]
Every neuron is included and the active sets may differ at the two times.
For the noise-scaled normal form it is
$\Lambda=\eta^2(kd-zc)$, not either cross entry separately.

\begin{proposition}[Minimal transverse snapshot certificate]
\label{sbend:leakage-certificate}
At two reached times $t_{\rm spec}\le t_a<t_b\le H$, suppose
$|\Lambda_a|,|\Lambda_b|\le\ell<k$, $|X_a|\le\alpha$, and
$\Lambda_a-\Lambda_b\ge d_\Lambda>0$. Then
\begin{equation}\label{sbend:leakage-budget}
 E_b-E_a\ge (k-\ell)d_\Lambda
                          -\frac{\alpha^2+s_0S_H}{2}.
\end{equation}
A strictly positive right side supplies $b_*$ in
Proposition~\ref{sbend:four-snapshots}.
\end{proposition}
\begin{proof}
Only the possibly negative longitudinal difference is bounded by
$-X_a^2/2$. The installed ambient contraction gives
$\|P_a\|\le\|W_\perp(t_a)\|_F\sqrt{S(t_a)}
\le\sqrt{s_0S_H}$. Discard the nonnegative final ambient square.
The remaining estimate is exactly \eqref{sbend:leakage-secants}.
\end{proof}

Thus removing a \emph{negative} coordinate-2 output that helped reconstruct
this probe can supply a rebound. In contrast, if $f_2\ge0$ and only this
coordinate decreases while the others remain fixed, the error decreases:
the derivative of its squared residual is $(k+f_2)\dot f_2\le0$.
Likewise ambient output has no target component here; reducing its norm
alone improves reconstruction. The contraction of the stacked ambient
weights does not by itself imply pointwise contraction of $P(\xi)$ because
input activations also move. These signs prevent using unspecified
``leakage decay'' as the missing later witness.

For example, a held positive strong family has $z_{i2}\ge0$. Its own
contribution to $\Lambda$ is nonpositive at every probe, irrespective of
the input gates. A nearly positive-$e_2$ weak family is inactive on this
negative-$e_2$ sector whenever its input angles exceed $\arcsin r$.
Beneficial negative transverse output, if used, must therefore be tracked
in the actual active complement or in a strong family not yet known to
have positive output coordinates. Positivity of the aggregate training
coefficient $c_{21}^O$ does not settle this probe-selected quantity.

\subsection{A zero-lag endpoint mechanism: removing a useful probe response}

Small lag is not an obstruction to the snapshot route. A perfectly
matched coherent planar block with unit direction $n(t)$ and mass $m(t)$
has $f_B(\xi)=m n(n^\top\xi)_+$ and, with no background,
\begin{equation}\label{sbend:matched-benefit}
 E(\xi,t)=\frac12-\frac12m(t)(2-m(t))(n(t)^\top\xi)_+^2.
\end{equation}
For $0<m\le1$, losing the positive probe activation can return the
error toward $1/2$ even with identically zero input--output angular lag.
The gate-exit derivative vanishes at the exact zero-lag exit, but the
finite endpoint gain in \eqref{sbend:matched-benefit} need not vanish.
This calculation is algebraic; a matched rotating block is not asserted
to follow the actual gradient flow.

\begin{proposition}[Post-specialization gate removal with signed background]
\label{sbend:sweep-certificate}
Fix a cohort $B$, two reached times $t_{\rm spec}\le t_a<t_b\le H$,
and a probe $\xi=(\sin\theta,-\cos\theta)$. Let $0<g\le1/4$.
At $t_a$ suppose every $i\in B$ has nonzero projected radii and
\[
 g\le\theta-\phi_i\le5g/4,\qquad
 g\le\theta-\psi_i\le5g/4,\qquad
 m_*\le M_B:=\sum_{i\in B}r_is_i\le1.
\]
At $t_b$ suppose all its probe gates are off. Let $G_j$ denote the full
output of $B^c$ at the probe at time $t_j$, let
$F_a=f_B(\xi,t_a)$, and define the signed background cost
\begin{equation}\label{sbend:background-cost}
 \mathfrak C_B:=\langle G_a,F_a\rangle
 +\frac12\|G_a-\xi\|^2-\frac12\|G_b-\xi\|^2
 +\frac12\|P_{\mathcal S^\perp}F_a\|^2.
\end{equation}
If $\mathfrak C_B\le7m_*g^2/96$, then
\begin{equation}\label{sbend:sweep-gap}
 E(\xi,t_b)-E(\xi,t_a)\ge m_*g^2/8>0.
\end{equation}
This conclusion requires no lower bound on angular lag. The background
contains every neuron outside $B$, including weak and unassigned neurons.
\end{proposition}
\begin{proof}
The two angle bounds give
$\xi^\top F_a\ge M_B\sin^2g$ and
$\|P_{\mathcal S}F_a\|\le M_B\sin(5g/4)\le5M_Bg/4$.
Since $M_B\le1$ and
$\sin^2g\ge g^2-g^4/3\ge47g^2/48$,
\[
 \xi^\top F_a-\frac12\|P_{\mathcal S}F_a\|^2
 \ge M_Bg^2\left(\frac{47}{48}-\frac{25}{32}\right)
 =\frac{19}{96}M_Bg^2.
\]
The exact endpoint expansion, using $f_B(\xi,t_b)=0$, is this
benefit minus $\mathfrak C_B$, leaving at least
$(19-7)m_*g^2/96=m_*g^2/8$.
\end{proof}

\subsection{Exact lag profile and its conditional viability}

Let $n_i=v_i/r_i$, $t_i=n_i^\angle$, $q_i=s_i/r_i$, and
$\zeta_i=\psi_i-\phi_i$. In this subsection $q_i$ is the radial ratio,
not the longitudinal ratio used in Section~\ref{sbpos:scope}.
For the actual compatible selected matrix write
\[
 a_i=n_i^\top A_i n_i,\quad b_i=n_i^\top A_i t_i,\quad
 c_i=t_i^\top A_i n_i,\quad d_i=t_i^\top A_i t_i,\quad
 p_i=t_i^\top P_i/r_i.
\]
The exact polar equations are
\[
 \dot\phi_i=q_i(b_i\cos\zeta_i+d_i\sin\zeta_i)-p_i,\qquad
 \dot\psi_i=q_i^{-1}(c_i\cos\zeta_i-a_i\sin\zeta_i).
\]
Thus, with
$\alpha_i=q_i^{-1}a_i+q_id_i$ and
$\beta_i=q_i^{-1}c_i-q_ib_i$,
\begin{equation}\label{sbend:lag-exact}
 \dot\zeta_i=-\alpha_i\sin\zeta_i+\beta_i\cos\zeta_i+p_i.
\end{equation}
There is no extra boundary approximation here: $A_i$ already uses the
actual compatible selectors. In particular
\begin{equation}\label{sbend:lag-skew}
 \beta_i=X_i+(q_i^{-1}-1)c_i-(q_i-1)b_i,\qquad
 X_i=\mathbb E_n[(f_1x_2-f_2x_1)\alpha_i(x)].
\end{equation}
For $q_i=1$ and zero ambient forcing, the forcing at zero lag is
exactly $X_i$, while the damping coefficient is $\operatorname{tr}A_i$.
Both depend on the learned residual; neither has an automatic late sign.

Put $y_i=\sin\zeta_i$, $k_i=\alpha_i\cos\zeta_i$, and
$F_i=\beta_i\cos^2\zeta_i+p_i\cos\zeta_i$. Then the exact, scalar
variation-of-constants identity is
\begin{equation}\label{sbend:lag-profile}
 y_i(t)=e^{-\int_a^t k_i}y_i(a)
             +\int_a^t e^{-\int_v^t k_i}F_i(v)\,dv.
\end{equation}
For example, on a stopped interval with
$0<\underline k\le k_i\le\overline k$ and $|F_i|\le F_*$,
\[
 |y_i(t)|\le e^{-\underline k(t-a)}|y_i(a)|
       +\frac{F_*}{\underline k}(1-e^{-\underline k(t-a)}).
\]
If instead $F_i\le-F_-<0$ there, the one-sided estimate is
\begin{equation}\label{sbend:lag-lower}
 -y_i(t)\ge
 \frac{F_-}{\overline k}(1-e^{-\overline k(t-a)})
                 -e^{-\underline k(t-a)}[y_i(a)]_+.
\end{equation}
Thus forcing of order $e$ can sustain lag of order $e$ over a fixed
window. Monotone lag decay is not a consequence of the mismatch damping
term. Conversely, a nonzero lower lag bound requires a signed forcing
estimate such as the one in \eqref{sbend:lag-lower}; an upper Cartesian
mismatch bound cannot supply it.

At the outward exit normal, the zero-background contribution is
$-r_is_i\dot\phi_i y_i$. The full contribution plus background is
\[
 -r_is_i\dot\phi_i y_i
 -r_i\dot\phi_i\langle G,w_i\rangle
 +\langle G-\xi,\dot G\rangle.
\]
The last two terms are signed. An $O(e)$ lag and $O(e)$ angular velocity
only give an $O(r_is_i e^2)$ individual exit contribution; they do not
bound the rest of the network. This is why the endpoint route remains
primary even after the lag audit.

\section{A learned-mass endpoint theorem in an extreme-separation region}
\label{sbsep:scope}

This section proves a restricted noisy many-neuron result from isotropic
initialization, with a learned strong cohort of mass at least $1/2$ before
the rebound. It does not prove learned weak-family specialization. The
region permits $\mu_2/\mu_1$ and the noise to be very small relative to the
initialization scale. Consequently this is a candidate one-concept
O2b-endpoint theorem, not a closure of the original two-family theorem or
of its fixed concept-ratio regime. Both restrictions are theorem-scope
choices; neither is evidence of an obstruction to the original goal.

\subsection{Directional concentration removes an artificial width conflict}

\begin{lemma}[Directional initial fluctuation bound]
\label{sbsep:directional-variance}
In Theorem~\ref{sbaudit:initial-descent}, set
$C_\xi=\mathbb E_n[|\xi^\top x|\|x\|]$. One may replace $C_x$ by
$C_\xi$ in the definition of $b_h$, while retaining $C_x$ in the residual
cost. If $C_\xi\le C_{\rm dir}$ on a deterministic arc, the uniform-grid
proof likewise permits replacing $\overline C$ by $C_{\rm dir}$ in its
width requirement, but not in its net radius or initialization cap.
\end{lemma}
\begin{proof}
Write $v_x=x/\|x\|$ and use weights
$|\xi^\top x|\|x\|/C_\xi$ in Jensen's inequality. Then
\[
 |Y_i|\le C_\xi,\qquad
 \mathbb EY_i^2\le C_\xi^2\mathbb E_n^{\rm weighted}
 \mathbb E[(\widehat u^\top\xi)^2(\widehat u^\top v_x)^2]
 \le\frac{3C_\xi^2}{d(d+2)}.
\]
The Bernstein and net arguments are unchanged. In particular a target
margin of order $r$ and $C_{\rm dir}=O(r)$ do not force $h\gtrsim r^{-2}$.
\end{proof}

\subsection{Exact all-neuron rank-one reference and uniform comparison}

Set $\mu_1=1$. The reference dataset has equal mass at $e_1$ and zero;
its time variable is $\tau=t/2$. Let $I_+=\{i:u_{i1}(0)>0\}$.
On these neurons write the input matrix as $(a^\top,B^\top)^\top$ and
the output matrix as $(c^\top,D^\top)^\top$, separating $e_1$ from its
orthogonal complement. Here $a,c\in\mathbb R^{|I_+|}$ and
$B,D\in\mathbb R^{(d-1)\times|I_+|}$. While the strong training gates
are retained, $B$ is constant and
\begin{equation}\label{sbsep:exact-system}
 a'=(1-a^\top c)c-D^\top Da,\qquad
 c'=(1-a^\top c)a,\qquad D'=-Daa^\top.
\end{equation}
All neurons in $I_+^c$ are stationary. Balanced isotropic initialization
gives $a(0)=c(0)=a_0$, $D(0)=B$. Therefore
\begin{equation}\label{sbsep:gram}
 cc^\top+D^\top D-aa^\top=B^\top B,\qquad
 a'=c-\|a\|^2a-B^\top Ba,\qquad \|D\|_F\le\|B\|_F.
\end{equation}
This reference includes every neuron, not a selected two-neuron subsystem.

Put $A_0=\|a_0\|$, $v=a_0/A_0$, and define
\begin{equation}\label{sbsep:scalar-reference}
 A'=A(1-A^2),\quad A(0)=A_0,\qquad
 d_0(\tau)=\sqrt{\frac{1-A(\tau)^2}{1-A_0^2}},\qquad
 D^\dagger=B[I-(1-d_0)vv^\top].
\end{equation}
The superscript $\dagger$ below means $a^\dagger=c^\dagger=Av$ and
$D^\dagger$ as above, with the actual constant $B$ and the actual inactive
neurons. It is an approximation to \eqref{sbsep:exact-system}, not an
assertion that all input or output directions coincide.

\begin{lemma}[Relative comparison through learned mass]
\label{sbsep:relative-comparison}
Suppose $\|B\|_F^2\le s_0$, $A_0^2\le s_0$, $v_i\ge v_*>0$,
and $0<\nu\le1/64$. On a scalar-reference interval $0\le\tau\le T\le H$
with $A^2\le1-\nu$, set
\[
 \rho=\frac{3s_0H}{\nu},\qquad
 \delta_c=4\rho+\frac{4s_0}{\nu},\qquad
 D_* =\frac{36s_0^{3/2}H^2}{\nu}.
\]
If
\begin{equation}\label{sbsep:comparison-caps}
 s_0\le\frac{\nu^2}{64(H+1)},\qquad \delta_c\le v_*/4,
\end{equation}
then all strong gates are retained throughout and
\begin{align}
 \|a-Av\|&\le4A\rho,&\quad\|c-Av\|&\le A\delta_c,
 &\quad\|D-D^\dagger\|_F&\le D_* .\label{sbsep:comparison-bounds}
\end{align}
The total reference mass is at most $1+2s_0$. If $A^2\ge3/4$, then
$a^\top c\ge2/3$ and $a_i,c_i\ge v_*/2$.
\end{lemma}
\begin{proof}
Let $u=\|a\|$, $p=c-a$, $l=\|p\|/u$ and
$\vartheta=a^\top p/u^2$. Stop before a strong gate changes or either
$u^2=1-\nu/2$ or $l=\nu/8$. Equations \eqref{sbsep:gram} give
\[
 \frac{u'}u=1-u^2+\vartheta-\widehat a^\top B^\top B\widehat a=:g,
 \qquad p'=-(1-a^\top c)p+D^\top Da.
\]
On the stopped interval $g\ge\nu/4$, and
$1-a^\top c+g\ge\nu/2$. In the upper-Dini sense,
$l'\le-(\nu/2)l+s_0$, hence $l\le2s_0/\nu$.
Writing $q=\log(u/A)$ yields
$q'=-A^2(e^{2q}-1)+\epsilon_q$, $|\epsilon_q|\le3s_0/\nu$.
The restoring sign gives $|q|\le3s_0\tau/\nu$ without an exponential
in $H$. Also $\|\widehat a'\|\le l+s_0\le3s_0/\nu$, so
$\|\widehat a-v\|\le3s_0\tau/\nu$.
These bounds give the first two inequalities in
\eqref{sbsep:comparison-bounds}. The caps make
$\rho<\nu/8$ and exclude both radial stopping faces:
$(1-\nu)e^{\nu/4}<1-\nu/2$ and $2s_0/\nu<\nu/8$.
They also give $a_i,c_i\ge3Av_i/4$, which excludes gate exit and
proves continuation from the initially positive gates.

The error $Z=D-D^\dagger$ satisfies
\[
 Z'=-Zaa^\top-D^\dagger(aa^\top-A^2vv^\top).
\]
The first term dissipates $\|Z\|_F^2$.
Since $\|aa^\top-A^2vv^\top\|\le12A^2(3s_0\tau/\nu)$,
integration gives the stated $D_*$ (with a factor-two reserve).
Finally $\|a\|^2<1$, $\|B\|_F^2\le s_0$, and the inactive input
mass is at most $s_0$. This gives the full mass bound. At $A^2\ge3/4$,
\[
 a^\top c\ge A^2e^{-2\rho}(1-2s_0/\nu)
 \ge\tfrac34(63/64)^2>2/3,
\]
and $a_i,c_i\ge3Av_i/4\ge v_*/2$.
\end{proof}

\subsection{The signed endpoint gain}

Let $B_2$ be the row of $B$ corresponding to $e_2$ and set
$b=Bv$, $b_2=B_2v$, $\eta=|b_2|>0$, and $s=\operatorname{sign}b_2$.
Choose $r=8\sqrt{s_0}/v_*\le1/8$, $k=\sqrt{1-r^2}$,
$\xi_0=re_1+sk e_2$, and $z=r/\eta$.
Assume $\eta\ge\eta_*>0$, let $Z_*=r/\eta_*$, and take
$\nu=1/(512Z_*^2)$. In particular $8\le z\le Z_*$.
Define scalar-reference times by
\begin{equation}\label{sbsep:snapshot-times}
 A(\tau_{\rm spec})^2=3/4,\qquad
 (A(\tau_a)^2-1)z+A(\tau_a)k=0,\qquad
 A(\tau_b)^2=1-\nu.
\end{equation}
Then $\tau_{\rm spec}<\tau_a<\tau_b$. These times may depend on
initialization; all will have a deterministic upper bound below.
For $\tau\ge\tau_{\rm spec}$ every $I_+$ probe gate is strictly on.
Indeed $ra_i\ge rv_*/2=4\sqrt{s_0}>|B_{2i}|$.

Write $G$ for the stationary output of $I_+^c$ at $\xi_0$, so
$\|G\|\le s_0$. Put $C=BB_2^\top-bb_2$, with $\|C\|\le2s_0$.
The approximate full output is exactly
\begin{equation}\label{sbsep:probe-reference}
 f^\dagger(\xi_0)=
 e_1(A^2r+A\eta k)+(0,skC)+\eta(0,b)T+G,
 \qquad T=d_0(Az+k).
\end{equation}
The transverse component along $b$ initially helps reconstruct the
probe and subsequently decreases; the inactive complement $G$ is
retained in the endpoint squares.

\begin{lemma}[Endpoint rebound at learned mass]
\label{sbsep:reference-rebound}
In addition to Lemma~\ref{sbsep:relative-comparison}, suppose
\begin{equation}\label{sbsep:rebound-caps}
 s_0\le\min\{1/256,1/(32Z_*^2)\},\quad
 s_0^{3/2}\le\eta_*/16,\quad s_0^2\le\eta_*^2/8,
\end{equation}
and
\begin{equation}\label{sbsep:output-cap}
 \epsilon_f:=4r\delta_c+3rD_*+4r\sqrt{s_0}\rho
 \le\eta_*^2/128.
\end{equation}
Then the exact rank-one reference obeys
$E(\xi_0,2\tau_b)-E(\xi_0,2\tau_a)\ge3\eta^2/8$.
Here $E(\xi_0,2\tau)$ denotes the reference error at physical time
$t=2\tau$.
\end{lemma}
\begin{proof}
The unique root defining $\tau_a$ has $A_a>9/10$: evaluate its quadratic
at $A=9/10$ and use $z\ge8$. Moreover
$1-A_a^2=A_ak/z>\nu$, proving the strict ordering.
At the two endpoints,
\[
 T_a\ge\tfrac9{16}\sqrt z,\qquad
 T_b\le\sqrt{2\nu}(Z_*+1)\le1/8,
 \qquad T_a-T_b\ge1,\quad 0\le T_a,T_b\le2Z_*.
\]
Set $Q=skC+G_\perp-sk e_2$, viewing $e_2$ in the $(d-1)$-dimensional
transverse space. Then
\[
 Q^\top b=-\eta k+(skC+G_\perp)^\top b,
 \qquad |(skC+G_\perp)^\top b|\le3s_0^{3/2}\le\eta k/4.
\]
Thus the transverse error secant in \eqref{sbsep:probe-reference} is
\[
 \eta(T_b-T_a)Q^\top b+
 \tfrac12\eta^2\|b\|^2(T_b^2-T_a^2)
 \ge\tfrac9{16}\eta^2-2\eta^2s_0Z_*^2
 \ge\tfrac8{16}\eta^2.
\]
At $\tau_a$ the longitudinal residual of the variable part is zero.
Its full initial longitudinal residual is $G_1$, so its secant is
bounded below by $-G_1^2/2\ge-\eta^2/16$.
The approximate full secant is therefore at least $7\eta^2/16$.

At either endpoint, \eqref{sbsep:comparison-bounds} and
$\|ra+skB_2^\top\|\le2r+\sqrt{s_0}$ give
$\|f-f^\dagger\|\le\epsilon_f$; the first-coordinate cost is at most
$\delta_c(3r+\sqrt{s_0})\le4r\delta_c$, and the transverse cost is
at most $3rD_*+4r\sqrt{s_0}\rho$.
The exact reference has mass at most $1+2s_0$, and $\epsilon_f<1/16$,
so an error perturbation of norm $\epsilon_f$ changes each error square
by at most $4\epsilon_f$. The two-endpoint cost is at most
$8\epsilon_f\le\eta^2/16$, as required.
\end{proof}

\subsection{Nonempty primitive noisy region and high-probability initialization}

The following explicit choices are deliberately conservative. They certify
nonemptiness, not a useful experimental parameter range. All constants
are primitive; none presupposes a trajectory event.
Fix $d\ge2$, $0<\delta<1/4$, and an even total sample count $N\ge2$.
Set
\[
 G_*=2+\lceil256\pi d\rceil,\qquad \ell=\log(8G_*/\delta),
\]
and choose the integer width so that
\begin{equation}\label{sbsep:width}
 h\ge\max\{1536(16\pi)^2\ell,
           (1024\pi/3)d\ell,32\log(16/\delta)\}.
\end{equation}
Define
\begin{align*}
 \beta&=\frac{\delta}{32h\sqrt d},&
 v_*&=\frac\beta{\sqrt h},&
 b_*&=\frac{\delta^3}{262144\,d\,h^{3/2}},\\
 s_0&=h\varepsilon^2,& \eta_*&=\varepsilon b_*,&
 r&=\frac{8\sqrt{s_0}}{v_*},\\
 Z_*&=\frac r{\eta_*}=\frac{8\sqrt h}{v_*b_*},&
 \nu&=\frac1{512Z_*^2},&
 H&=\frac12\log\frac4{\nu s_0\beta^2}.
\end{align*}
Use $\rho,\delta_c,D_*,\epsilon_f$ from the preceding lemmas with
this $H$. Choose $\varepsilon>0$ sufficiently small that all of
\eqref{sbsep:comparison-caps}, \eqref{sbsep:rebound-caps}, and
\eqref{sbsep:output-cap} hold, together with
\begin{equation}\label{sbsep:remaining-caps}
 r\le1/8,\qquad s_0\le\frac r{384\pi d},\qquad
 \frac{8\sqrt{s_0}\rho}{v_*^2}
       +\frac{3(\delta_c+D_*)}{v_*}\le\eta_*/40.
\end{equation}
Finally set
\begin{align}
 \chi_*&=\min\{1/32,\varepsilon\beta/8,
                   \eta_*^2/1024,\eta_*v_*/240\},\notag\\
 q_D&=\chi_*e^{-20H},\label{sbsep:data-caps}\\
 0<\mu_2^2&\le\min\{r^2/100,q_D/256\},\notag\\
 0<a_D&\le\min\{\mu_2/4,r^2/100,
                           \varepsilon\beta/16,q_D/160\},\notag\\
 0<\sigma_1,\sigma_2&\le
                   a_D/\sqrt{2\log(4N/\delta)}.\notag
\end{align}
Here $\sigma_k$ are standard deviations of planar isotropic Gaussian
noise, there are $N/2$ independent samples per cluster, and $\mu_1=1$.

\begin{theorem}[Noisy many-neuron strong-cohort specialization and endpoint rebound]
\label{sbsep:theorem}
For the primitive region above, balanced isotropic independent
initialization $u_i(0)=w_i(0)=\varepsilon\widehat u_i$ has the following
properties with probability at least $1-\delta$ jointly with the samples.
There is a fixed initialization cohort $I_+$ of at least $h/4$ neurons,
and times
\[
 0<t_{\rm early}<t_{\rm spec}<t_a<t_b\le2H,
\]
such that throughout $[t_{\rm spec},t_b]$ its input and output vectors
lie within angle $r/2$ of $e_1$ and its learned longitudinal mass satisfies
\begin{equation}\label{sbsep:learned-mass}
 \sum_{i\in I_+}u_{i1}(t)w_{i1}(t)\ge1/2.
\end{equation}
The entire network has $S(t)\le4$ on $[0,t_b]$.
On each of the two deterministic arcs
\[
 \mathcal I_\pm=\{\sin\theta\,e_1\pm\cos\theta\,e_2:
       |\theta-\arcsin r|\le r/4\},\qquad
 \Delta_{\rm init}=r/2,
\]
there is uniform initial decrease with the explicit
$\tau_{\rm dec},c_{\rm dec},t_{\rm early},a_{\rm dec}$ of
Theorem~\ref{sbend:descent-clock}, taking
$J_*=r/(8\pi)$ and $\overline C=\overline\lambda=2$.
At one of their centers $\xi_0$ the later secant is
\begin{equation}\label{sbsep:noisy-rebound}
 E(\xi_0,t_b)-E(\xi_0,t_a)\ge b_{\rm reb}:=\eta_*^2/8>0.
\end{equation}
On an open sector, the initial drop is at least $a_{\rm dec}/2$.
The rebound after specialization is at least $b_{\rm reb}/2$.
An explicit lower bound on its angular width is
\begin{equation}\label{sbsep:final-sector}
 \min\{r/4,4\arcsin(\rho_{\rm swing}/2)\},\qquad
 \rho_{\rm swing}=\min\{1,\min(a_{\rm dec},b_{\rm reb})/100\}.
\end{equation}
Every such error has an interior minimum on $[0,t_b]$.
The conclusion does not locate the minimum after $t_{\rm spec}$, fit the
weak concept, or create a second learned family.
\end{theorem}
\begin{proof}
\emph{The initialization geometry.}
Let $R_i^2=\widehat u_{i1}^2+\widehat u_{i2}^2$ and
$R_*^2=\delta/(32dh)$. With failure probability less than $\delta/4$,
\begin{equation}\label{sbsep:geometry-event}
 |I_+|\ge h/4,\quad
 \min_i|\widehat u_{i1}|\ge\beta,\quad
 \min_iR_i^2\ge R_*^2,\quad |B_2v|\ge\varepsilon b_*.
\end{equation}
Indeed the first failure is at most $e^{-h/8}\le\delta/32$.
The density of one spherical coordinate on $[-\beta,\beta]$ is at most
$\sqrt d$, giving a union bound $2h\sqrt d\beta=\delta/16$.
For $d>2$, $R_i^2$ has distribution function
$1-(1-q)^{(d-2)/2}\le dq$ for $0\le q\le1/2$; for $d=2$ it is one.
The radius failure is at most $\delta/32$.
Condition on all radii, signs of the first coordinates, and all but the
first positive neuron's planar angle. Its angle is uniform on a
half-circle. Its contribution to
$\sum_{I_+}\widehat u_{i1}\widehat u_{i2}$ is
$(R_i^2/2)\sin(2\theta_i)$. The elementary arcsine small-ball bound
\[
 \sup_z\Pr\{|\sin(2\theta_i)-z|\le u\}\le4\sqrt u
\]
follows by integrating its density on an interval, with the largest
endpoint singularity proportional to $1/\sqrt{1-|x|}$.
Taking $q=R_*^2\delta^2/8192$ bounds the probability that this sum has
absolute value less than $q$ by $\delta/16$.
Since $A_0\le\varepsilon\sqrt h$, this proves the last bound in
\eqref{sbsep:geometry-event}. This conditioning is on signs and radii,
not on the minimum-coordinate event; a union bound combines them.
Moreover
\[
 s_0\beta^2/4\le A_0^2\le s_0,\qquad v_i\ge\beta/\sqrt h=v_*.
\]
The explicit logistic solution gives
$\tau_b=\tfrac12\log[(1-\nu)(1-A_0^2)/(\nu A_0^2)]\le H$.
Apply the preceding reference lemmas up to this time.

\emph{Uniform noisy comparison, including arbitrary weak gates.}
On the sample event $\max\|x^{(k)}-\mu_ke_k\|\le a_D$, the failure
probability is at most $\delta/4$ by the planar Gaussian radial tail
and a union bound over $N$ samples. Couple the real and reference
networks with the same initialization. Stop while both full input and
output Frobenius norms are at most $2$, and all strong-sample gates
agree with $I_+$. The reference vector field, in physical time and
combined parameter Frobenius norm, is $10$-Lipschitz there.
For example its longitudinal-input and output difference bounds have
coefficients $2$ and $9/2$, whose block norm is at most $13/2<10$.
At a fixed parameter point the strong-sample covariance changes by at
most $a_D+a_D^2/2$. Since $\|I-WU_{I_+}^\top\|\le5$, the combined
strong-gradient cost is at most $30a_D$. The entire weak gradient,
with its actual arbitrary selectors, has combined norm at most
$10(\mu_2+a_D)^2$. Thus a sufficient forcing bound is
\[
 P_D=40[a_D+(\mu_2+a_D)^2].
\]
The parameter distance up to physical time $2H$ is at most
\begin{equation}\label{sbsep:noisy-distance}
 \chi:=P_De^{20H}\le\chi_*.
\end{equation}
This is a finite-window comparison with the vanishing data perturbation,
not a Gronwall estimate on the $B^\top B$ amplification error; that error
was controlled relatively in Lemma~\ref{sbsep:relative-comparison}.
The caps in \eqref{sbsep:data-caps} give $P_D<q_D$.

Reference strong margins are at least $\varepsilon\beta/2$ for $I_+$,
and at least $\varepsilon\beta$ in absolute value for $I_+^c$.
The perturbations cost at most $\chi+2a_D\le\varepsilon\beta/4$.
Reference parameter norms are at most $\sqrt{1+2s_0}$, so adding
$\chi\le1/32$ also excludes the norm stopping face. The comparison
therefore continues to $2\tau_b\le2H$. In particular $S\le4$
through the last snapshot. The deterministic $2H$ is used only as an
upper bound in the comparison integral; no radial estimate past
$2\tau_b$ is needed.

At the reached times $t_j=2\tau_j$, the learned strong mass changes
by at most $4\chi$ from its reference value $2/3$, leaving at least
$1/2$. Its longitudinal coordinates are at least $v_*/2-\chi$ and its
transverse input/output norms at most $\sqrt{s_0}+\chi$, giving angles
at most $4\sqrt{s_0}/v_*=r/2$. At either rebound snapshot the full
output changes by at most $4\chi$, hence its error changes by at most
$16\chi$. The secant loses at most $32\chi\le\eta_*^2/32$, which is
less than the reserve in Lemma~\ref{sbsep:reference-rebound}.
This proves \eqref{sbsep:noisy-rebound}.
The limitation concerning the weak concept is quantitative: throughout this
window, transverse input and output Frobenius norms are each at most
$\sqrt{s_0}+\chi$, so
$|f_2(e_2,t)|\le(\sqrt{s_0}+\chi)^2\le2s_0$.
By homogeneity this is also the coordinate-2 fitting fraction at
$\mu_2e_2$. Thus the theorem does not hide a second learned family.

\emph{The initialized decreasing event on the same primitive region.}
Uniformly on both arcs, the first coordinate of the probe is in
$[3r/4,5r/4]$. The bounded-noise event and the data caps give
$C_x\le2$, $C_\xi\le2r$ and $J_\xi\ge r/(8\pi)$.
For the last assertion use $cF(c)\ge c$ on cluster 1 and the crude
negative bound $cF(c)\ge-1$ on cluster 2. Its numerator is at least
\[
 [(3r/4)(1-a_D)-a_D](1-a_D)-(\mu_2+a_D)^2\ge r/4.
\]
This estimate holds for both signs of the second probe coordinate.
A net with chord mesh $r/(512\pi d)$ requires at most $G_*$ points
per arc. Lemma~\ref{sbsep:directional-variance}, the width condition,
and a union bound over both nets give failure probability at most
$\delta/4$. Keep $\overline C=2$ for grid transport and the initial
residual cost. The initialization cap in
\eqref{sbsep:remaining-caps} is precisely the required one.
Theorem~\ref{sbend:descent-clock} now gives the claimed deterministic
initial drop. Since $A_0^2\le1/256$, $2\tau_{\rm spec}>1$ whereas
$t_{\rm early}<1$, proving the strict ordering.
The geometry, data and descent failures sum to less than $\delta$;
no independence between their derived events is assumed.

Finally apply the four-snapshot argument with the proven mass cap $4$
at the four snapshots. It gives \eqref{sbsep:final-sector}.
All unit probes in these sectors are at distance at least $3/4$ from
every training sample under the displayed data caps, so the statement
concerns extrapolation probes even if the selected sign is positive.

\emph{Nonemptiness.}
Choose $d,\delta,N,h$ first. The constants $v_*,b_*,Z_*,\nu$ do not
depend on $\varepsilon$. As $\varepsilon\downarrow0$,
$H=O(\log(1/\varepsilon))$, $\rho,\delta_c=O(\varepsilon^2H)$,
$D_*=O(\varepsilon^3H^2)$ and
$\epsilon_f=O(\varepsilon^3H)$, whereas $\eta_*^2$ is a positive
constant times $\varepsilon^2$. Every initialization cap is therefore
satisfied for sufficiently small positive $\varepsilon$. Then choose
positive $\mu_2,a_D,\sigma_1,\sigma_2$ successively from
\eqref{sbsep:data-caps}. The region is nonempty and all time, probability
and endpoint budgets hold there simultaneously.
\end{proof}

\subsection{The lag remains of initialization scale after learned mass}

This proof also identifies the lag profile; it does not assume damping
has removed it. For $i\in I_+$ and $\tau_{\rm spec}\le\tau\le\tau_b$,
put $q_i=B_{2i}/(Av_i)$. In the approximation,
\begin{equation}\label{sbsep:lag-profile}
 \zeta_i^\dagger(\tau)=
 \arctan\left(q_i-\frac{(1-d_0)b_2}{A}\right)-\arctan q_i.
\end{equation}
Here $|q_i|\le r/4$, $d_0\le\sqrt{(1/4)/(1-s_0)}$, and
$A\ge\sqrt{3/4}$. The mean-value theorem gives
\[
 \eta/5\le-s\zeta_i^\dagger\le2\eta.
\]
The reference comparison and the noisy parameter perturbation give
\[
 |\zeta_i-\zeta_i^\dagger|
 \le\frac{8\sqrt{s_0}\rho}{v_*^2}
       +\frac{3(\delta_c+D_*)}{v_*}+\frac{6\chi}{v_*}
 \le\eta_*/20.
\]
Consequently the actual planar lag obeys the uniform learned-window bound
\begin{equation}\label{sbsep:lag-lower}
 \eta/8\le-s\zeta_i(t)\le3\eta
 \qquad(t_{\rm spec}\le t\le t_b).
\end{equation}
Thus this trajectory retains a nonzero lag of order $|B_2v|$, inherited
from initialization, after learning order-one strong mass. It is not a
noise-generated lag. No gate-exit witness is used: all strong probe gates
are on during the two endpoint snapshots. The decaying signed transverse
coefficient in \eqref{sbsep:probe-reference} supplies the rebound.

The scope question left by this result is exact: whether an initialized
noisy many-neuron theorem with one learned strong cohort and
initialization-dependent extreme concept separation is an acceptable
narrowing, or whether the paper requires two learned families and/or a
concept ratio bounded away from zero. The theorem above establishes the
former mathematical result; it does not choose it as the paper's headline.

\end{document}
